\SourceLeaf{137}
\SixthArticle{6}
Die vorangehenden Funktionen sind algebraische Funktionen sowie transzendente Funktionen, die von den Logarithmen und dem Kreis abhängen. Keineswegs jedoch unternehmen wir unsere \textit{allgemeine} Untersuchung um \textit{dieser} willen, sondern vielmehr, um die Theorie der höheren transzendenten Funktionen voranzubringen, von denen unsere Reihe eine sehr umfangreiche Klasse umfasst. Dazu gehören unter unendlich vielen anderen die Koeffizienten, die entstehen, wenn die Funktion $(aa+bb-2ab\cos\varphi)^{-n}$ in eine Reihe nach den Kosinus der Winkel $\varphi$, $2\varphi$, $3\varphi$ usw. entwickelt wird. Über diese werden wir bei anderer Gelegenheit \textit{im Besonderen} ausführlicher sprechen. Diese Koeffizienten lassen sich auf mehrere Arten auf die Form unserer Reihe zurückführen. Setzen wir nämlich
\[
(aa+bb-2ab\cos\varphi)^{-n}=\Omega=A+2A'\cos\varphi+2A''\cos2\varphi+2A'''\cos3\varphi+\mbox{usw.}
\]
so haben wir \textit{erstens}
\[
\begin{array}{rl}
A&=a^{-2n}F(n,n,1,\frac{bb}{aa})
\end{array}
\]
\[
\begin{array}{rl}
A'&=na^{-2n-1}bF(n,n+1,2,\frac{bb}{aa})
\end{array}
\]
\[
\begin{array}{rl}
A''&=\frac{n(n+1)}{1\mathbin{.}2}a^{-2n-2}bbF(n,n+2,3,\frac{bb}{aa})
\end{array}
\]
\[
\begin{array}{rl}
A'''&=\frac{n(n+1)(n+2)}{1\mathbin{.}2\mathbin{.}3}a^{-2n-3}b^3F(n,n+3,4.\ \frac{bb}{aa})
\end{array}
\]
\begin{center}usw.\end{center}
Betrachtet man nämlich $aa+bb-2ab\cos\varphi$ als Produkt aus $a-br$ und $a-br^{-1}$ (wobei $r$ die Größe $\cos\varphi+\sin\varphi.\BareRoot{-1}$ bezeichnet), so wird $\Omega$ gleich dem Produkt
\[
\begin{array}{rl}
\mbox{aus}&a^{-2n}
\end{array}
\]
\[
\begin{array}{rl}
\mbox{und}&1+n\frac{br}{a}+\frac{n(n+1)}{1\mathbin{.}2}.\ \frac{bbrr}{aa}+\frac{n(n+1)(n+2)}{1\mathbin{.}2\mathbin{.}3}.\ \frac{b^3r^3}{a^3}+\mbox{usw.}
\end{array}
\]
\[
\begin{array}{rl}
\mbox{und}&1+n\frac{br^{-1}}{a}+\frac{n(n+1)}{1\mathbin{.}2}.\ \frac{bbr^{-2}}{aa}+\frac{n(n+1)(n+2)}{1\mathbin{.}2\mathbin{.}3}.\ \frac{b^3r^{-3}}{a^3}+\mbox{usw.}
\end{array}
\]
Da dieses Produkt mit
\[
A+A'(r+r^{-1})+A''(rr+r^{-2})+A'''(r^3+r^{-3})
\]
identisch sein muss, ergeben sich die oben angegebenen Werte unmittelbar.

Ferner haben wir \textit{zweitens}
\SourceLeaf{138}
\[
\begin{array}{rl}
A&=(aa+bb)^{-n}F(\frac12 n,\frac12 n+\frac12 .\,1,\frac{4aabb}{(aa+bb)^2})
\end{array}
\]
\[
\begin{array}{rl}
A'&=n(aa+bb)^{-n-1}abF(\frac12 n+\frac12,\frac12 n+1,2,\frac{4aabb}{(aa+bb)^2})
\end{array}
\]
\[
\begin{array}{rl}
A''&=\frac{n(n+1)}{1\mathbin{.}2}(aa+bb)^{-n-2}aabbF(\frac12 n+1,\frac12 n+\frac32,3,\frac{4aabb}{(aa+bb)^2})
\end{array}
\]
\[
\begin{array}{rl}
A'''&=\frac{n(n+1)(n+2)}{1\mathbin{.}2\mathbin{.}3}(aa+bb)^{-n-3}a^3b^3F(\frac12 n+\frac32,\frac12 n+2,4,\frac{4aabb}{(aa+bb)^2})
\end{array}
\]
\begin{center}usw.\end{center}
Diese Werte ergeben sich leicht aus
\[
\Omega(aa+bb)^n=1+n(r+r^{-1})\frac{ab}{aa+bb}+\frac{n(n+1)}{1\mathbin{.}2}(r+r^{-1})^2\frac{aabb}{(aa+bb)^2}+\mathrm{usw.}
\]
\Needspace{10\baselineskip}
\textit{Drittens} gilt
\[
\begin{array}{rl}
A&=(a+b)^{-2n}F(n,\frac12,1,\frac{4ab}{(a+b)^2})
\end{array}
\]
\[
\begin{array}{rl}
A'&=n(a+b)^{-2n-2}abF(n+1,\frac32,3,\frac{4ab}{(a+b)^2})
\end{array}
\]
\[
\begin{array}{rl}
A''&=\frac{n(n+1)}{1\mathbin{.}2}(a+b)^{-2n-4}aabbF(n+2,\frac52,5,\frac{4ab}{(a+b)^2})
\end{array}
\]
\[
\begin{array}{rl}
A'''&=\frac{n(n+1)(n+2)}{1\mathbin{.}2\mathbin{.}3}(a+b)^{-2n-6}a^3b^3F(n+3,\frac72,7,\frac{4ab}{(a+b)^2})
\end{array}
\]
\begin{center}usw.\end{center}
Schließlich gilt \textit{viertens}
\[
\begin{array}{rl}
A&=(a-b)^{-2n}F(n,\frac12,1,-\frac{4ab}{(a-b)^2})
\end{array}
\]
\[
\begin{array}{rl}
A'&=n(a-b)^{-2n-2}abF(n+1,\frac32,3,-\frac{4ab}{(a-b)^2})
\end{array}
\]
\[
\begin{array}{rl}
A''&=\frac{n(n+1)}{1\mathbin{.}2}(a-b)^{-2n-4}aabbF(n+2,\frac52,5,-\frac{4ab}{(a-b)^2})
\end{array}
\]
\[
\begin{array}{rl}
A'''&=\frac{n(n+1)(n+2)}{1\mathbin{.}2\mathbin{.}3}(a-b)^{2n-6}a^3b^3F(n+3,\frac72,7,-\frac{4ab}{(a-b)^2})
\end{array}
\]
\begin{center}usw.\end{center}
Jene und diese Werte ergeben sich leicht aus
\[
\begin{array}{rl}
\Omega(a+b)^{2n}&=(1-\frac{4ab\cos\frac12\varphi^2}{(a+b)^2})^{-n}
\end{array}
\]
\[
\begin{array}{rl}
&=1+n\frac{ab}{(a+b)^2}(r^{\frac12}+r^{-\frac12})^2+\frac{n(n+1)}{1\mathbin{.}2}\mathbin{.}\frac{aabb}{(a+b)^4}(r^{\frac12}+r^{-\frac12})^4+\mathrm{usw.}
\end{array}
\]
\[
\begin{array}{rl}
\Omega(a-b)^{2n}&=(1+\frac{4ab\sin\frac12\varphi^2}{(a-b)^2})^{-n}
\end{array}
\]
\[
\begin{array}{rl}
&=1+n\frac{ab}{(a-b)^2}(r^{\frac12}-r^{-\frac12})^2+\frac{n(n+1)}{1\mathbin{.}2}\mathbin{.}\frac{aabb}{(a-b)^4}(r^{\frac12}-r^{-\frac12})^4+\mathrm{usw.}
\end{array}
\]
\begin{center}\rule{.23\linewidth}{.4pt}\end{center}
